%@ pages 233-241
%@ status translated
% Original section 21, "Exemple: la descente pour les complexes", pp. 233--241.

\clearpage
\section{Example: descent for complexes}
\label{complexepage}

When one speaks of stacks, the example of complexes has always been present
both in algebraic geometry and in algebraic topology
(see \cite{SGA4b}, \cite{SGA6}, \cite{GabrielZisman}, \cite{Quillen}, \cite{ResiduesDuality},
\cite{Illusie}, \cite{ThomasonTrobaugh},
\cite{BondalKapranov}, \cite{Hinich2}, \cite{HoveyBook} for example).
It is moreover this example of complexes that drew us into this work, and
that guided us in particular in the formulation and proof of the results
of \S17, \S18 and \S19.


Let $\Xx$ be a site equipped with a sheaf of rings $\Oo$.
We assume that $\Xx$ admits fibre products and has sufficiently many disjoint
sums compatible with fibre products. For $X\in \Xx$, let
$Cpx_{\Oo}(X)$ be the category of complexes of sheaves of
$\Oo$-modules on $\Xx /X$. Let
$$
qis(X)\subset Cpx_{\Oo}(X)
$$
be the subcategory of arrows which are quasi-isomorphisms.

Let us denote by $Cpx_{\Oo}^{[0, \infty )}(X)$ the subcategory of complexes
concentrated in degrees
$\geq 0$. By
Quillen \cite{Quillen}, $Cpx_{\Oo}^{[0,
\infty )}(X)$ has a structure of closed model category in which the
weak equivalences are the quasi-isomorphisms and the cofibrations
are the injections.

Recently Hinich \cite{Hinich2} and Hovey \cite{HoveyBook} have equipped
the category $Cpx_{\Oo}(X)$ of all complexes (not necessarily
bounded) with a cmc structure whose
weak equivalences are the quasi-isomorphisms and whose cofibrations
are the injections.

In fact, Quillen as well as Hinich (and Hovey in an earlier version
of \cite{HoveyBook})
work in the dual situation, i.e.\ where the fibrations are the
surjections; one must therefore
apply their constructions
to the abelian category dual to that of $\Oo$-modules. This
duality could
possibly
raise problems of a set-theoretic nature. Hovey alludes to this in
his book
\cite{HoveyBook}, and attributes to Grodal
the construction in which the cofibrations are
the injections.
In the most recent version of \cite{HoveyBook}, Hovey gives the
construction of this cmc, which he
calls the ``injective cmc structure for complexes'', and he proves
that it is generated
by cofibrations
(\cite{HoveyBook} Def. 2.3.13). The reader may also recover this result by applying the criterion
\ref{dhklemme}
to the category $Cpx_{\Oo}(X)$ with $W = qis(X)$ and $cof$ the term-by-term
injections of complexes. One verifies properties (4) and (5) of
\ref{dhklemme}
using once again Jardine's technique \cite{Jardine} (the other
properties are
easy).


% much of the following discussion was changed in v3

We therefore indeed have a cmc $Cpx_{\Oo}(X)$ generated by cofibrations, and one
verifies that the functor
$$
X\mapsto Cpx_{\Oo}(X)
$$
is a left Quillen presheaf on $\Xx$.

We would like to apply
\ref{dansleschamps} in order to obtain that this presheaf is a stack.
Unfortunately, in the present version {\tt v3} we cannot
do this, because we do not know how to compute directly the
simplicial presheaves of morphisms for unbounded complexes.

We therefore restrict ourselves to the case of complexes bounded below
(we thus withdraw our statement for unbounded complexes).

Let $Cpx_{\Oo}^{\{ 0,\infty )}(X)$ be the subcategory of complexes
concentrated
cohomologically in degrees $\geq 0$.
 By \ref{stabilite}, $L(Cpx_{\Oo}^{\{ 0,\infty )}(X))\subset
L(Cpx_{\Oo}(X))$ is a full simplicial subcategory.
Descent preserves the condition of cohomological truncation. An argument
based on truncation of complexes shows that
$$
L(Cpx_{\Oo}^{[0,\infty )}(X))\rightarrow L(Cpx_{\Oo}^{\{ 0,\infty )}(X))
$$
is an equivalence. One can equip $Cpx_{\Oo}^{[0,\infty )}(X)$ with a
structure of cmc generated by cofibrations (this is easier than Hovey's
result mentioned above for unbounded complexes).


Now we apply Theorem
\ref{dansleschamps} to obtain the following corollary,
namely that the localised category $L(Cpx_{\Oo}^{[0,\infty )})$
is a Segal
$1$-stack.

It should be noted
that this corollary is
essentially the ``cohomological descent'' for complexes of SGA 4
\cite{SGA4b} Exposé Vbis,
but with an additional component,
``strictification of descent data''
(see \S18), which does not appear in SGA 4. For this corollary,
one could probably replace the part of
the proof of
\ref{dansleschamps} which is found in \S19 by the arguments of \cite{SGA4b}.

%%% note that the statement of this corollary was changed in v3;
%%% we now look only at bounded complexes

\begin{corollary}
\label{cestunchamp}
Suppose that $\Xx$ admits fibre products and sufficiently many disjoint
sums. Then the presheaf of simplicial categories
$$
L(Cpx_{\Oo}^{[0,\infty )}):=
\left( X\mapsto  L(Cpx_{\Oo}^{[0,\infty )}(X), qis(X)) \right)
$$
is a Segal $1$-stack on $\Xx$, which we call the {\em Segal $1$-stack of
modules of complexes on $(X,\Oo )$}.
\end{corollary}
{\em Proof:} Apply \ref{dansleschamps}. Compatibility (3) with
products is immediate, as is compatibility (2) with disjoint
sums. Condition (1) is a consequence of the local character of the
notion of quasi-isomorphism. For condition (0) one applies the theorem
\ref{sectionscmf}
(II), saying as usual that the verification of the first
part of the hypothesis
(the one corresponding to conditions (4) and (5) of \ref{dhklemme}) calls on
the techniques of
\cite{Jardine}.

%%% verification added in v2:
%%% this is page 229 in v1

For condition (4) (added in version 2 of the paper) we shall prove
that if $f: X\rightarrow Y$ is a morphism in $\Xx $ then the restriction
functor $f^{\ast}$ for complexes, which is left Quillen, is also right
Quillen. This assertion immediately implies the desired compatibility
with limits. In order to prove that $f^{\ast}$
is right Quillen one constructs its left adjoint $f_{!}$ by
analogy with the construction of \S4. Since we are dealing with complexes and not
with $n$-categories, one replaces the disjoint union by the direct sum.
For a complex $C^{\cdot}$ of sheaves of $\Oo$-modules
on $\Xx /X$, one sets:
$$
f_{!, 0}(C^{\cdot})(Z) := \bigoplus _{\varphi : Z\rightarrow X}
C^{\cdot} (Z).
$$
Then one sets $f_!(C^{\cdot})$ equal to the sheaf associated to
$f_{!,0}(C^{\cdot})$. (If the representable objects are sheaves
on $\Xx$ this step is not necessary, because $f_{!,0}(C^{\cdot})$
is already a complex of sheaves.)

This functor $f_!$ is the left adjoint of $f^{\ast}$ and
the only problem is
to prove that it preserves cofibrations (resp.\ trivial cofibrations).
The preservation of cofibrations is immediate, because the cofibrations are
just
the injections (the functor $f_{!,0}$ preserves injections trivially
and the ``associated sheaf'' functor sends injections of presheaves
to injections of sheaves).

For trivial cofibrations it suffices to
prove that the functor $f_{!}$ preserves quasi-isomorphisms.
One observes that the functor
$f_{!,0}$ preserves the exact structure of complexes of presheaves,
and hence
that the cohomology sheaves of $f_{!}(C^{\cdot})$ are the sheaves
associated to the images under $f_{!,0}$ of the cohomology presheaves of
$C^{\cdot}$. Using the exact structure, it suffices to prove that if $A$
is a presheaf on $\Xx /X$ whose associated sheaf is $0$,
then the same holds for $f_{!,0}(A)$. The hypothesis on $A$ means
that every section becomes trivial on a covering. A section $a$ of
$f_{!,0}(A)$ is a finite sum of sections of $A$, so one can take
a common refinement of the coverings which trivialise these components,
in order to trivialise $a$. This completes the proof of condition (4),
which allows us to apply \ref{dansleschamps}.


%%%% end of the v2 addition

We now turn to the verification of condition (5)
of \ref{dansleschamps}.
We begin by identifying the simplicial presheaves of morphisms
in $L(Cpx_{\Oo}^{[0, \infty )})$. Let $DP$ be the
Dold--Puppe construction (\cite{DoldPuppe} cf.\ \cite{Illusie}) which
transforms a
presheaf of complexes of abelian groups into a presheaf
of simplicial abelian
groups; and let us denote by $\tau ^{\leq 0} C$ the ``smart'' truncation
(i.e.\ the one
which preserves cohomology) in degrees $\leq 0$ of a complex $C$. For
$U,V\in Ob\,Cpx_{\Oo}^{[0, \infty )}(X)$, one has
$$
L(Cpx_{\Oo}^{[0, \infty )})_{1/}(U,V)
\cong DP(\tau ^{\leq 0} {\bf R} \underline{Hom} (U,V))
$$
where, for $Y\in \Xx /X$,
$$
{\bf R} \underline{Hom} (U,V)(Y):=
Hom _{\Oo}(U|_{\Xx /Y}, V'|_{\Xx /Y})
$$
denotes
the {\em internal $\underline{Hom}$ presheaf}, with $V\rightarrow V'$ a
quasi-isomorphism to a complex of injective sheaves.
More generally, if $\Ff$ is an abelian category with sufficiently
many injectives, if $Cpx^{[0, \infty )}(\Ff )$ is the category of complexes
in degrees $\geq 0$ of objects of $\Ff$, then one obtains the simplicial
category $L(Cpx^{[0, \infty )}(\Ff ), qis)$. One has the same formula
$$
L(Cpx^{[0, \infty )}(\Ff ), qis)_{1/}(U,V)
\cong DP(\tau ^{\leq 0} {\bf R} \underline{Hom} (U,V)).
$$
In particular, $L(Cpx^{[0, \infty )}(\Ff ), qis)$ is very close
to the {\em differential graded category} of complexes considered
by Bondal--Kapranov \cite{BondalKapranov}. On the other hand the truncation
$$
\tau _{\leq 2} L(Cpx^{[0, \infty )}(\Ff ), qis),
$$
which can be chosen as
a strict $2$-category, is the one which was considered by
Gabriel--Zisman \cite{GabrielZisman}.

Let us return to our complexes of sheaves.
We shall prove directly condition (5) of \ref{dansleschamps}
(in other words part (a) of criterion \ref{critere}), namely that
$$
L(Cpx_{\Oo}^{[0, \infty )})_{1/}(U,V)
$$
is a stack. Let
$$
I^{\cdot}:= {\bf R} \underline{Hom} (U,V)
:= \underline{Hom} (U,V');
$$
this is a complex
of sheaves of acyclic abelian groups. For every $k\geq 0$ one
defines the truncation $U_{\leq k}$ of the complex $U$ (this
truncation preserves
cohomology). This truncation is a subcomplex of $U$ and $U$ is
the colimit of the $U_{\leq k}$. The morphisms in the system are
injections i.e.\ cofibrations of complexes, so the colimit is an
$hocolim$. One sets
$$
I^{\cdot}_{k} :=\underline{Hom} (U_{\leq k},V').
$$
Then $I^{\cdot }_{k}$ is a complex concentrated in degrees $\geq -k$
whose components are sheaves of acyclic abelian
groups.

It follows that
$DP(\tau ^{\leq 0} I^{\cdot}_{ k})$ admits a (finite) dévissage by
successive homotopy fibre products
with simplicial presheaves
of the form $DP(I^j[j]) = K(I^j, j)$, where the $I^k$ are acyclic;
the latter are stacks
by \cite{Jardine} \cite{KBrown}, and at the start ($j=0$)
one has $K(\ker (I^0_k\rightarrow I^1_k), 0)$, which is a sheaf
of sets, hence a stack. Applying Lemma \ref{fiprodchamp} one obtains
that $DP(\tau ^{\leq 0} I^{\cdot}_{k})$ is a stack.


On the other hand, $I^{\cdot} = \lim _{\leftarrow , k} I^{\cdot }_{k}$.
One obtains that
$$
DP(\tau ^{\leq 0} I^{\cdot}) = \lim _{\leftarrow , k}
DP(\tau ^{\leq 0} I^{\cdot}_k)
$$
and this limit is an object-by-object homotopy limit over
$\Xx$:
the transition morphisms
are surjections of simplicial abelian groups
(except in degree zero). In other words
the morphisms in the inverse system are HBKQ fibrations
of simplicial presheaves. If one replaces
this inverse system by a system, object-by-object equivalent,
denoted $\{ A_k\}$, whose objects are fibrant and whose transition
morphisms are
fibrations for the Jardine structure, then for every object
$X$ of $\Xx$ one will have
$$
DP(\tau ^{\leq 0} I^{\cdot}) (X) = \lim _{\leftarrow , k}
DP(\tau ^{\leq 0} I^{\cdot}_k)(X)
\cong \lim _{\leftarrow , k} A_k(X).
$$
On the other hand each $A_k$ is a stack for the topology $\Gg$ of the site
$\Xx$, so by the lemmas
\ref {grofibchampimplfib} and
\ref{stuff} the $A_k$ are $\Gg$-fibrant
and the transition morphisms are $\Gg$-fibrations.
Hence the limit $\lim _{\leftarrow , k} A_k$ is a stack,
which proves that
$$
L(Cpx_{\Oo}^{[0, \infty )})_{1/}(U,V)
= DP(\tau ^{\leq 0} I^{\cdot})
$$
is a stack.
This completes the verification of condition (5) of the theorem
\ref{dansleschamps}, which allows us to conclude that
$L(Cpx_{\Oo}^{[0,\infty )})$ is a Segal $1$-stack.
\eop

We conjecture of course that the hypothesis on the site $\Xx$
is not
necessary. On the other hand, the weakening (in {\tt v3}) of this corollary,
due to the error
in Lemma \ref{pourloc}, prompts us to ask the following question:
\newline
---{\em Question:} Is the Dwyer--Kan localised Segal $1$-prestack
of unbounded complexes with respect to quasi-isomorphisms,
$L(Cpx_{\Oo})$, a stack?


We now pass to the case of complexes bounded
on both sides.
%%%%%%%%%%%%%%%%%%% end of the v3 changes %%%%%%%%%%

For $0\leq a < \infty$ let us denote by
$Cpx_{\Oo}^{[0, a\} }(X)$ (resp.\ $Cpx_{\Oo}^{[0, a]}(X)$)
the subcategory of complexes with vanishing cohomology
in degree $>a$ (resp.\ with
vanishing components in degree $>a$). Using property
\ref{stabilite} one obtains that
$$
L(Cpx_{\Oo}^{[0, a\} }(X), qis(X))\rightarrow
L(Cpx_{\Oo}^{[0, \infty )}(X), qis(X))
$$
is a full simplicial subcategory, and with an argument of truncation
of complexes (using for example \ref{stabilite}) one obtains that the
morphism
$$
L(Cpx_{\Oo}^{[0, a] }(X), qis(X))\rightarrow
L(Cpx_{\Oo}^{[0, a\} }(X), qis(X))
$$
is an equivalence. One obtains an object-by-object equivalence
between presheaves of simplicial categories on $\Xx$
$$
L(Cpx_{\Oo}^{[0, a] })\cong
L(Cpx_{\Oo}^{[0, a\} }).
$$
These are full Segal $1$-subprestacks of
$L(Cpx_{\Oo}^{[0, \infty )})$; in particular
condition (a) of \ref{critere} is
automatic; but since the vanishing of the cohomology in
degrees $>a$ is
a local property, condition (b) of \ref{critere} is verified by
$L(Cpx_{\Oo}^{[0, a\} })$. One concludes that
$L(Cpx_{\Oo}^{[0, a] })$ is a Segal $1$-stack.

By translation, $L(Cpx _{\Oo}^{[a,b]})$ is a Segal $1$-stack. For
$n=b-a$ the stacks of morphisms here are $n$-truncated simplicial
presheaves. In particular (cf.\ the notation of \S2),
$$
L(Cpx _{\Oo}^{[a,b]})\cong \Re _{\geq 1}
\Pi _n \circ L(Cpx _{\Oo}^{[a,b]}),
$$
so by abuse of notation one can identify $L(Cpx _{\Oo}^{[a,b]})$
with the $n+1$-stack (not a Segal one)
$$
\tau _{\leq n+1}L(Cpx _{\Oo}^{[a,b]})=\Pi _n \circ L(Cpx _{\Oo}^{[a,b]}).
$$
This $n+1$-stack is $1$-groupic (indeed it comes from an
$n+1$-truncated Segal $1$-stack).


\subnumero{Perfect complexes}

The main references are SGA 6 (\cite{SGA6}), Illusie \cite{Illusie} and
Thomason--Trobaugh \cite{ThomasonTrobaugh}, but there are many other
works on this subject.

Let $CpxParf^{[a,b]}(X)$ be the category of perfect complexes of
$tor$-amplitude contained in the interval $[a,b]$, and let $W$ be the
subcategory of quasi-isomorphisms. By the theory of Dwyer--Kan
\cite{DwyerKan3}, cf.\ \S8, one obtains $\infty$-categories
$L( CpxParf^{[a,b]}(X))$
which are in fact
strict simplicial categories and which we regard as
Segal $1$-categories, but which are on the other hand $n+1$-truncated for
$n=b-a$ and which one could therefore regard as $n+1$-categories. When
$X$ varies this provides an $n+1$-prestack ${\cal A}^{[a,b]}$
over $\Xx$:
$$
{\cal A}^{[a,b]} (X):= \Pi _n \circ L( CpxParf^{[a,b]}(X)).
$$

As a corollary of \ref{cestunchamp} one has:

\begin{proposition}
\label{intro1}
The $n+1$-prestack ${\cal  A}^{[a,b]}$ is an $n+1$-stack ($n=b-a$).
\end{proposition}
\eop

Let $CpxFib^{[a,b]}(X)$ be the category of complexes of bundles
(i.e.\ of locally free sheaves of $\Oo$-modules of finite rank) with support
in the interval $[a,b]$. One defines a new $n+1$-prestack
$L(CpxFib^{[a,b]})$ by setting
$$
L(CpxFib^{[a,b]})(X):= L(CpxFib^{[a,b]}(X)).
$$
This $n+1$-prestack comes with a morphism
$$
L(CpxFib^{[a,b]})\rightarrow {\cal  A}^{[a,b]}=L(CpxParf^{[a,b]}).
$$

One introduces the subcategory $CpxProj^{[a,b]}(X)\subset CpxFib^{[a,b]}(X)$
of complexes of projective $\Oo (X)$-modules. Note that to every complex
of projective $\Oo (X)$-modules $M$ one can associate the complex of sheaves
on $\Xx/X$
$$
\tilde{M}:Y\mapsto M_Y:= M\otimes _{\Oo
(X)}\Oo (Y).
$$
On $CpxProj^{[a,b]}(X)$ one has at one's disposal Quillen's
cmc structure, in which one has
$$
L(CpxProj^{[a,b]}(X))_{1/}(U,V)
\cong DP(\tau ^{\leq 0} {\bf R} \underline{Hom} (U,V)).
$$
What is involved here is the ${\bf R} \underline{Hom}$ of complexes of
$\Oo (X)$-modules. In general it differs from the
${\bf R} \underline{Hom}$ of sheaves, i.e.\ the arrow
$$
{\bf R} \underline{Hom} (U,V) \rightarrow {\bf R} \underline{Hom}
(\tilde{U},\tilde{V}) (X)
$$
is not in general an equivalence. This is due to the non-triviality
in cohomology of $X$ (in the site of schemes one could restrict
to affine $X$ and would not have this problem). However one can see
(let us admit that more details would have to be supplied on this point,
which will be done in a later work) that the morphism of prestacks on $\Xx /X$:
$$
\left( Y\mapsto {\bf R} \underline{Hom} _{\Oo (Y)}(U_Y,V_Y)\right)
\rightarrow {\bf R} \underline{Hom}
(\tilde{U},\tilde{V})
$$
is a $\Gg$-weak equivalence. Now one has
$$
L(CpxParf^{[a,b]})_{1/}(\tilde{U}, \tilde{V})=
DP \tau ^{\leq 0}{\bf R} \underline{Hom}
(\tilde{U},\tilde{V})
$$
and in particular, the latter is a stack. Hence it is the stack associated
to the prestack
$$
L(CpxProj^{[a,b]})_{1/}(U,V) =
\left( Y\mapsto {\bf R} \underline{Hom} _{\Oo (Y)}(U_Y,V_Y)\right) .
$$
By Lemma \ref{formule} one concludes that the morphism
$$
({\bf ch}L(CpxProj^{[a,b]}))_{1/}(U,V) \rightarrow
L(CpxParf^{[a,b]})_{1/}(\tilde{U}, \tilde{V})
$$
is an equivalence, i.e.\ that
$$
{\bf ch}L(CpxProj^{[a,b]})\rightarrow L(CpxParf^{[a,b]})
$$
is fully faithful. The morphism
$$
L(CpxProj^{[a,b]})\rightarrow L(CpxParf^{[a,b]})
$$
is clearly $\Gg$-essentially surjective, so one obtains that this
morphism induces an equivalence
$$
{\bf ch}L(CpxProj^{[a,b]})\stackrel{\cong}{\rightarrow} L(CpxParf^{[a,b]}).
$$
It follows
(using the chain of inclusions $CpxProj \subset CpxFib \subset CpxParf$
and the result of the preceding line) that ${\bf ch}L(CpxFib^{[a,b]})$
retracts onto ${\bf ch}L(CpxProj^{[a,b]})$.

One sees on the other hand, from the description of $L(CpxFib^{[a,b]})(X)$ by
``hammocks'' (Dwyer--Kan \cite{DwyerKan2}, cf.\ \S8), that every $i$-arrow in
this $n+1$-category comes locally (with respect to the topology
$\Gg$) from an $i$-arrow of $L(CpxProj^{[a,b]})$. This implies that the
retraction found is an equivalence
$$
{\bf ch}L(CpxProj^{[a,b]}) \cong {\bf ch}L(CpxFib^{[a,b]}).
$$
We have finally obtained the following proposition.


\begin{proposition}
${\cal  A}^{[a,b]}=L(CpxParf^{[a,b]})$ is the $n+1$-stack associated to the
$n+1$-prestack $L(CpxFib^{[a,b]})$.
\end{proposition}
\eop

We now indicate an alternative proof of Illusie's result
\cite{Illusie}
on the existence of the symmetric powers etc.\ of perfect complexes. This
proof
gives life to the intuition according to which ``a perfect complex
being locally a
complex of vector bundles,
it suffices to take the symmetric power of each of these
complexes of bundles and to
glue them together''.

\begin{corollary}
Any construction (such
as the symmetric power) which comes from a functor
(natural in $X\in \Xx$)
$$
CpxFib^{[a,b]}(X) \rightarrow CpxFib^{[a',b']}(X)
$$
and which is compatible with weak equivalences,
extends to a morphism of $n+1$-stacks
$$
L(CpxParf^{[a,b]})\rightarrow L(CpxParf^{[a',b']}).
$$
This morphism truncates in turn to a morphism between the derived
categories
$$
{\cal D}^{[a,b]}_{\rm parf}(X)\rightarrow {\cal D}^{[a',b']}_{\rm parf}(X).
$$
\end{corollary}
{\em Proof:} Apply the universal property of the stack associated to a
prestack (see \S13) to obtain the morphism between the associated
stacks. By
noting the equality
$$
{\cal D}^{[a,b]}_{\rm parf}(X)=\tau _{\leq 1} \left(
L(CpxParf^{[a,b]})(X)\right) ,
$$
one obtains the morphism between the derived categories.
\eop



We now turn to the case of the site $\Xx = Sch$ of schemes.
One then has at one's disposal a notion of {\em locally geometric
$n$-stack} of
$n$-groupoids \cite{geometricN}, which is the analogue for
$n$-stacks of the
notion of Artin's {\em algebraic $1$-stack} \cite{ArtinInventiones}. The
notion extends to $n$-stacks not necessarily of groupoids. Without
entering into the details of this definition, one can say that if $A$ is
a (locally) geometric $n$-stack then its interior $A^{int,
0}$ is
a (locally) geometric $n$-stack of groupoids in the sense of
\cite{geometricN}. We state here a result which will be the subject of another
work.

\begin{theorem}
\label{modparfait}
On the site of schemes,
the $n+1$-stack ${\cal  A}^{[a,b]}$ is locally geometric ($n=b-a$).
It is covered by the open substacks ${\cal  A}^{s_a,\ldots ,s_b}$
obtained by imposing on the cohomology at every point the bound $h^i \leq
s_i$;
these substacks are geometric.
\end{theorem}

In fact, the reader may perhaps find the proof of this
theorem by relying on the ``homological perturbation
theory'' of Gugenheim {\em et al.} \cite{Gugenheim}
\cite{GugenheimLambe}. What is essentially involved is to prove that the
morphism from the evident parameter scheme for a complex of trivial bundles
(of given ranks) to the stack ${\cal  A}^{[a,b]}$ is smooth.



% The original opens the parenthesis "(dont les valeurs sont ... (X)" and
% never closes it; we print the closing parenthesis after $(X)$.
The $1$-truncation $\tau _{\leq 1}{\cal  A}^{[a,b]}$ (whose values
are the $1$-truncations of the $n+1$-categories ${\cal  A}^{[a,b]}(X)$)
is the presheaf of categories which to $X$ associates the {\em derived
category} of perfect complexes of amplitude contained in $[a,b]$. Neither
the theorem \ref{intro1} nor
the theorem \ref{modparfait} remains true for this
truncation: this shows the interest of ``considering higher
homotopies''.

One can use the $L(Cpx _{\Oo})$ and ${\cal A}^{[a,b]}$ to define
the notions of complex,
perfect complex etc.\ on Segal $n$-stacks (in particular, on
$1$-stacks, e.g.\ Artin's
algebraic $1$-stacks). If $B$ is a Segal $n$-stack, we
shall say that a {\em complex of
$\Oo$-modules on $B$} is a morphism (to the fibrant replacement)
$$
B\rightarrow L(Cpx
_{\Oo})'.
$$
A {\em perfect complex of amplitude contained in $[a,b]$} is a morphism
$$
B\rightarrow ({\cal A}^{[a,b]})'.
$$
The results above show that if $B$ is the $0$-stack represented
by an object of the site
(e.g.\ a scheme),
then these definitions coincide with the usual definitions.

\newpage
